\section*{Streszczenie}
\label{chap:Streszczenie}

Celem pracy było zaprojektowanie i~implementacja systemu RiskScoop AI -- prototypu agenta AI do analizy ryzyka powodziowego dla infrastruktury krytycznej. System łączy przetwarzanie języka naturalnego z~analizą geoprzestrzenną, umożliwiając zadawanie pytań dotyczących zagrożenia powodziowego.

Praca przedstawia przegląd literatury obejmujący agentów AI w~analizie geoprzestrzennej oraz systemy Rapid Mapping. Zidentyfikowano lukę rynkową -- brak narzędzi łączących interfejs języka naturalnego z~analizą ryzyka powodziowego.

System integruje dane z~Overture Maps Foundation (infrastruktura), Copernicus GloFAS (zagrożenie powodziowe) oraz model Gemini do interpretacji zapytań. Aplikacja webowa (Mapbox GL JS) wizualizuje wyniki na interaktywnej mapie z~kwantyfikacją ryzyka.

Prototyp poddano testom funkcjonalnym i~wydajnościowym. System obsługuje zapytania dotyczące kategorii infrastruktury (szpitale, szkoły, drogi), automatycznie tłumacząc je na operacje geoprzestrzenne. Wyniki prezentowane są na mapie z~kodowaniem kolorystycznym poziomu ryzyka.

Praca ma zastosowanie w~obszarze Rapid Mapping oraz wsparcia decyzyjnego dla służb reagowania kryzysowego.

\bigskip

\textbf{Słowa kluczowe}: \quad agent AI \quad analiza ryzyka powodziowego \quad interfejs języka naturalnego \quad infrastruktura krytyczna \quad Rapid Mapping

\textbf{Keywords}: \quad AI agent \quad flood risk analysis \quad natural language interface \quad critical infrastructure \quad Rapid Mapping
